\documentclass[10pt,a4paper]{amsart}
\usepackage[margin=19mm]{geometry}
\usepackage{amsmath,amssymb,mathtools}
\usepackage[scr=rsfs,cal=euler]{mathalfa}
\usepackage{xfrac}
\usepackage[unicode,hidelinks,bookmarks=false]{hyperref}
\newcommand{\ie}{i.e.\ }
\newcommand{\cf}{cf.\ }
\newcommand{\resp}{respectively\ }
\newcommand{\Subsection}{\subsection}
\providecommand{\nobreakdash}{\nobreak\hbox{-}}
\setlength{\emergencystretch}{2em}
\multlinegap0pt
\usepackage[shortlabels]{enumitem}
\usepackage{amssymb}
\usepackage{mathtools}
\newtheorem{restatable}{Restatable}
\newtheorem{setup}{Set-up}
\newtheorem{proposition}{Proposition}
\newtheorem{thm}{Theorem}
\usepackage{cleveref}
\crefname{restatable}{Restatable}{Restatables}
\crefname{setup}{Set-up}{Set-ups}
\crefname{proposition}{Proposition}{Propositions}
\crefname{thm}{Theorem}{Theorems}
\newcommand{\C}{\mathbb{C}}
\renewcommand{\O}{\mathcal{O}}
\renewcommand{\P}{\mathbb{P}}
\newcommand{\tC}{\widetilde{C}}
\title{Low Degree Points on Singular Plane Curves}
\author{Zachary R. Canale, Nathan Chen, Zoe Curewitz, Jacob A. Daum, Karina Dovgodko, Carlos F. Santiago-Calder\'on,, Shiv R. Yajnik}
\date{Source: arXiv:2603.20529v1 (CC BY 4.0)}
\begin{document}
\maketitle

\noindent\emph{Source and licence.} Low Degree Points on Singular Plane Curves. Version arXiv:2603.20529v1 (CC BY 4.0), distributed by arXiv under the Creative Commons Attribution 4.0 International licence, http://creativecommons.org/licenses/by/4.0/, as recorded in the arXiv licence metadata for this exact version and shown on the abstract page https://arxiv.org/abs/2603.20529v1. Full licence notice: \texttt{licenses/arxiv-2603.20529-CC-BY-4.0.txt}.\par\smallskip

\noindent\textit{Editorial scope.} Counted unit: the article's thm CCCDDSY2 (source0.tex lines 411--413). The article's complete original proof of it is delivered with it (source0.tex lines 414--423, one \texttt{proof} environment of 10 lines). No mathematics is written, repaired or re-derived: the statement and the proof are the article's own lines, in the article's own notation, and no retained line is substituted or re-flowed.  Retained uncounted as the source-local context the proof needs: lines 129--132; lines 155--157; lines 180--187; lines 405--407.  Nothing was omitted from any retained span, so the package carries no omission record.  No retained line contains a citation, so the wrapper carries no bibliography and no bibliography entry.  No retained line contains a figure environment, an image-inclusion command or drawing code, so no image is delivered and the package declares no asset dependency.  Editorial formatting only, in addition: an AMS \texttt{amsart} preamble that restates the article's own declarations and macros used by the retained lines, namely the environments \texttt{restatable}, \texttt{setup}, \texttt{proposition}, \texttt{thm} on separate counters, together with the standard packages those lines need.  The retained environments print in this document as Restatable 1, Setup 1, Proposition 1, Theorem 1 in order of appearance, whereas the article prints the same items with the numbering of its own full document.  The complete native source of the article as distributed in the arXiv e-print for this exact version is bundled unchanged as \texttt{sources/196855/source0.tex}.
\begin{restatable}{thm}{CCCDDSY}\label{thm:CCCDDSY} Let $C \subset \P^{2}$ be a plane curve over a number field $K$ with $\delta$ ordinary nodes and/or cusps. Suppose one of the following conditions holds:
\[ \{ d \geq 8 \quad \text{and} \quad \delta \leq d-8 \} \quad \text{or} \quad \left\{ d \geq 17 \quad \text{and} \quad \delta < \frac{d^{2}}{4} - 2d - \frac{17}{4} \right\}. \]
Then $C$ only contains finitely many points of degree $\leq d-3$, and all but finitely many points of $C$ of degree $\leq d-1$ arise as the intersection of $C$ with lines in $\P^2$ defined over $K$.
\end{restatable}

\begin{setup}\label{Setup}
    Fix a number field $K$ and let $C \subset \P^{2}$ be an integral plane curve over $K$ of degree $d$ with $\delta$ ordinary nodes and/or cusps (and no other singularities). Let $\nu \colon \tC \rightarrow C$ denote the normalization.
\end{setup}

\begin{proposition}\label{prop:corollaryOfCK}
    Let $C$ and $\tC$ be as in Set-up~\ref{Setup}.
    \begin{enumerate}[topsep=-1ex, itemsep=-0.5ex, partopsep=1ex, parsep=1ex]
        \item[(i)]\label{prop:corollaryOfCK-1} If $d \geq 4$ and $\delta < \frac{1}{4} d^{2} - d + \frac{11}{4}$, then $\tC$ does not admit a $g^{1}_{e}$ with $e \leq d-3$.
        \item[(ii)]\label{prop:corollaryOfCK-2} If $d \geq 5$ and $\delta \leq \frac{d^{2}}{4} - d$, then every $g^{1}_{d-2}$ on $\tC$ is induced by a pencil of lines.
        \item[(iii)]\label{prop:corollaryOfCK-3} If $d \geq 9$ and $\delta \leq \frac{d^{2}}{4} - d - 1$, then every $g^{1}_{d-1}$ on $\tC$ is induced by a pencil of lines.
    \end{enumerate}
\end{proposition}

\begin{thm}\label{thm: main_geometric}
Let $S/\C$ be a nice surface with $h^1(S, \O_S)=0$, and let $C$ be a smooth curve on $S$ whose class is big and nef. Then for $e < C^2/9$, the locus $W_e(C)$ contains no positive-dimensional abelian varieties.
\end{thm}

\begin{restatable}{thm} {CCCDDSY2}\label{thm:CCCDDSY2} Let $C \subset \P^{2}$ be a plane curve over a number field $K$ with $\delta$ ordinary nodes/cusps. Suppose $d \geq 9$ and $0<\delta \leq \frac{d^2}{4}-\frac{9}{4}d+2$.
Then $C$ only contains finitely many points of degree $\leq d-3$, and all but finitely many points of $C$ of degree $\leq d-1$ arise as the intersection of $C$ with lines in $\P^2$ defined over $K$.
\end{restatable}

\begin{proof}
    Let $\tC$ be the strict transform of $C$ inside the blowup $S \rightarrow \mathbb{P}^2$ of $\P^{2}$ at the singular points of $C$. Then $S$ is a nice surface with $H^1(S,\mathcal{O}_S)=0$, and $\tC$ is smooth since $C$ only has ordinary nodes/cusps. Furthermore, one can check that $\tC^{2} > 0$ and $\tC$ must intersect non-negatively against any other curves on $S$. By the Nakai-Moishezon criterion it follows that $\tC$ is nef, and $\tC^{2} > 0$ implies that $\tC$ is also big. If $d-1 < \frac{\tC^2}{9}$, then Theorem~\ref{thm: main_geometric} implies that $W_e(\tC)$ contains no positive-dimensional abelian varieties for $e \leq d-1$. The same argument as in the proof of \cref{thm:CCCDDSY} %(using \ref{prop:corollaryOfCK})%
    then gives us the desired result.
    
    All that remains to check is that the inequality below holds:
    \begin{equation*}
        \frac{\tC^2}{9} = \frac{d^2-4\delta}{9} > d-1 \iff d^2-9d+9-4\delta >0.
    \end{equation*}
    Note that the assumption $d \geq 9$ is used to show that $\frac{d^2}{4}-\frac{9}{4}d+2 \geq 1$.
\end{proof}

\end{document}
